\chapter{Limit Theorems}
\label{ch:limits}

Probability is most useful when many random contributions are combined. Two
facts then emerge: averages settle down to a constant (the law of large numbers),
and the fluctuations around that constant take a universal bell shape (the
central limit theorem). To state them we must first say what it means for a
sequence of random variables to converge. This chapter defines the four standard
modes of convergence, proves how they relate, proves the Borel--Cantelli lemmas,
and then proves the laws of large numbers and the central limit theorem.

\section{Modes of convergence}

Let $X,X_1,X_2,\dots$ be random variables on the same probability space (except
for convergence in distribution, where the spaces may differ).

\begin{definition}[Four modes]
\begin{enumerate}[label=(\alph*)]
\item $X_n\to X$ \emph{almost surely} (a.s.) if $\Prob\bigl(\lim_nX_n=X\bigr)=1$.
\item $X_n\to X$ \emph{in probability}, written $X_n\xrightarrow{p}X$, if
$\Prob(\abs{X_n-X}>\epsilon)\to0$ for every $\epsilon>0$.
\item $X_n\to X$ \emph{in $L^p$} ($p\ge1$) if $\Exp\abs{X_n-X}^p\to0$.
\item $X_n\to X$ \emph{in distribution}, written $X_n\xrightarrow{d}X$, if
$F_{X_n}(x)\to F_X(x)$ at every $x$ where $F_X$ is continuous.
\end{enumerate}
\end{definition}

\begin{nuance}{why only continuity points?}
Let $X_n=1/n$ (deterministic) and $X=0$. Then $F_{X_n}(0)=0$ for all $n$, but
$F_X(0)=1$. We would not want to say that $X_n$ fails to converge to $0$, so we
ignore the jump point $x=0$. The restriction to continuity points of the \emph{limit}
fixes exactly this.
\end{nuance}

\begin{theorem}[Relations between modes]\label{thm:modes}
\begin{enumerate}[label=(\alph*)]
\item If $X_n\to X$ a.s., then $X_n\xrightarrow{p}X$.
\item If $X_n\to X$ in $L^p$, then $X_n\xrightarrow{p}X$.
\item If $X_n\xrightarrow{p}X$, then $X_n\xrightarrow{d}X$.
\item If $X_n\xrightarrow{d}c$ for a constant $c$, then $X_n\xrightarrow{p}c$.
\end{enumerate}
No other implication holds in general (see the examples below).
\end{theorem}
\begin{proof}
(a) Fix $\epsilon>0$ and let $A_n=\{\abs{X_m-X}\le\epsilon\text{ for all }m\ge n\}$. The
$A_n$ increase. If $X_n(\omega)\to X(\omega)$ then $\omega\in A_n$ for some $n$,
so $\Prob(\bigcup_nA_n)\ge\Prob(X_n\to X)=1$. By continuity of probability,
$\Prob(A_n)\to1$. Since $\{\abs{X_n-X}>\epsilon\}\subseteq A_n^c$, we get
$\Prob(\abs{X_n-X}>\epsilon)\to0$.

(b) Markov's inequality: $\Prob(\abs{X_n-X}>\epsilon)\le\Exp\abs{X_n-X}^p/\epsilon^p\to0$.

(c) Let $x$ be a continuity point of $F=F_X$ and $\epsilon>0$. Then
\[
F_{X_n}(x)=\Prob(X_n\le x)\le\Prob(X\le x+\epsilon)+\Prob(\abs{X_n-X}>\epsilon),
\]
because $X_n\le x$ and $\abs{X_n-X}\le\epsilon$ imply $X\le x+\epsilon$. Similarly
$F_{X_n}(x)\ge F(x-\epsilon)-\Prob(\abs{X_n-X}>\epsilon)$. Letting $n\to\infty$:
$F(x-\epsilon)\le\liminf F_{X_n}(x)\le\limsup F_{X_n}(x)\le F(x+\epsilon)$. Let
$\epsilon\downarrow0$ and use continuity of $F$ at $x$.

(d) The CDF of the constant $c$ is $0$ below $c$, $1$ at and above $c$, continuous
everywhere except $c$. For $\epsilon>0$ the points $c\pm\epsilon$ are continuity
points, so
$\Prob(\abs{X_n-c}>\epsilon)\le F_{X_n}(c-\epsilon)+1-F_{X_n}(c+\epsilon)\to0+1-1=0$.
\end{proof}

\begin{example}[The converses fail]\label{ex:conversefail}
\begin{enumerate}[label=(\roman*)]
\item \emph{In probability but not almost surely} (the ``typewriter sequence''). On
$\Omega=[0,1]$ with Lebesgue measure, let $X_n=\indic_{I_n}$ where $I_n$ runs through
the intervals $[0,1]$; $[0,\frac12],[\frac12,1]$; $[0,\frac13],[\frac13,\frac23],[\frac23,1]$;
and so on, sweeping across $[0,1]$ with shorter and shorter intervals. Then
$\Prob(X_n\ne0)=\abs{I_n}\to0$, so $X_n\xrightarrow{p}0$, yet every $\omega$ is in
infinitely many $I_n$, so $X_n(\omega)=1$ infinitely often and $X_n$ converges for no
$\omega$.
\item \emph{Almost surely but not in $L^1$.} Let $U\sim U(0,1)$ and
$X_n=n\indic_{\{U<1/n\}}$. For every $\omega$ with $U(\omega)>0$, $X_n(\omega)=0$ once
$n>1/U(\omega)$, so $X_n\to0$ a.s.; but $\Exp[X_n]=n\cdot\frac1n=1\not\to0$. Mass
escapes to infinity on a set of vanishing probability.
\item \emph{In distribution but not in probability.} Let $X\sim N(0,1)$ and
$X_n=-X$. All $X_n$ have the same law as $X$, so $X_n\xrightarrow{d}X$, but
$\abs{X_n-X}=2\abs X$ does not tend to $0$ in probability.
\end{enumerate}
\end{example}

\section{The Borel--Cantelli lemmas}

For events $A_1,A_2,\dots$ define
\[
\{A_n\text{ i.o.}\}=\bigcap_{n=1}^\infty\bigcup_{m=n}^\infty A_m,
\]
the event that infinitely many of the $A_n$ occur (``i.o.'' is ``infinitely often'').
(An outcome is in this set iff for every $n$ it lies in some $A_m$ with $m\ge n$.)

\begin{theorem}[First Borel--Cantelli lemma]\label{thm:bc1}
If $\sum_n\Prob(A_n)<\infty$, then $\Prob(A_n\text{ i.o.})=0$.
\end{theorem}
\begin{proof}
For each $n$, $\{A_k\text{ i.o.}\}\subseteq\bigcup_{m\ge n}A_m$, so by countable
subadditivity $\Prob(A_k\text{ i.o.})\le\sum_{m\ge n}\Prob(A_m)$, which is the tail of
a convergent series and tends to $0$ as $n\to\infty$.
\end{proof}

\begin{theorem}[Second Borel--Cantelli lemma]\label{thm:bc2}
If the events $A_n$ are \emph{independent} and $\sum_n\Prob(A_n)=\infty$, then
$\Prob(A_n\text{ i.o.})=1$.
\end{theorem}
\begin{proof}
It suffices to show $\Prob\bigl(\bigcap_{m\ge n}A_m^c\bigr)=0$ for every $n$, since
then the complement of $\{A_k\text{ i.o.}\}$, which is the union over $n$ of these
events, is a countable union of null events. By independence (of the complements,
Proposition~\ref{prop:indcomp}) and $1-x\le e^{-x}$, for $N>n$,
\[
\Prob\Bigl(\bigcap_{m=n}^NA_m^c\Bigr)=\prod_{m=n}^N\bigl(1-\Prob(A_m)\bigr)
\le\exp\Bigl(-\sum_{m=n}^N\Prob(A_m)\Bigr)\xrightarrow[N\to\infty]{}0,
\]
because the series diverges. Let $N\to\infty$ using continuity from above.
\end{proof}

\begin{nuance}{independence is essential in the second lemma}
Let $U\sim U(0,1)$ and $A_n=\{U<1/n\}$. Then $\sum\Prob(A_n)=\sum1/n=\infty$, but the
events are nested, and $\Prob(A_n\text{ i.o.})=\Prob(U=0)=0$. The two lemmas together
form a ``zero--one law'' for independent events: the probability of infinitely many
occurrences is $0$ or $1$, according to whether $\sum\Prob(A_n)$ converges or
diverges.
\end{nuance}

\begin{example}[Monkey at a typewriter]
A monkey types independently, each key uniformly from $27$ keys. Let $A_k$ be the
event that the $k$th block of $5$ consecutive keystrokes spells a fixed word (such as
``HELLO''). The $A_k$ are independent with probability $27^{-5}>0$ each, so
$\sum\Prob(A_k)=\infty$ and by the second lemma the word appears infinitely often
almost surely.
\end{example}

\begin{theorem}[Criteria for almost sure convergence]\label{thm:subseq}
\begin{enumerate}[label=(\roman*)]
\item If $\sum_n\Prob(\abs{X_n-X}>\epsilon)<\infty$ for every $\epsilon>0$, then
$X_n\to X$ almost surely.
\item If $X_n\xrightarrow{p}X$, then there is a subsequence $X_{n_k}\to X$ almost
surely.
\end{enumerate}
\end{theorem}
\begin{proof}
(i) For $\epsilon=1/j$, Theorem~\ref{thm:bc1} gives $\Prob(\abs{X_n-X}>1/j\text{ i.o.})=0$.
The union over $j\in\N$ of these null events is null, and outside it $\abs{X_n-X}\le1/j$
eventually for every $j$, i.e.\ $X_n\to X$.
(ii) Choose $n_1<n_2<\dots$ with $\Prob(\abs{X_{n_k}-X}>2^{-k})<2^{-k}$. By
Theorem~\ref{thm:bc1}, almost surely $\abs{X_{n_k}-X}\le2^{-k}$ for all large $k$.
\end{proof}

\section{Laws of large numbers}

Let $X_1,X_2,\dots$ be i.i.d.\ with mean $\mu$ and put $S_n=X_1+\dots+X_n$.

\begin{theorem}[Weak law of large numbers]\label{thm:wlln}
If $\Var(X_1)=\sigma^2<\infty$, then $S_n/n\xrightarrow{p}\mu$; in fact
$\Prob(\abs{S_n/n-\mu}>\epsilon)\le\sigma^2/(n\epsilon^2)$.
\end{theorem}
\begin{proof}
$\Exp[S_n/n]=\mu$ and $\Var(S_n/n)=n\sigma^2/n^2=\sigma^2/n$ by independence. Apply
Chebyshev's inequality.
\end{proof}

The strong law says more: the averages converge for almost every outcome. We prove it
under a finite variance using a maximal inequality.

\begin{theorem}[Kolmogorov's maximal inequality]\label{thm:kolmax}
Let $X_1,\dots,X_n$ be independent with $\Exp[X_i]=0$ and finite variances. Then for
every $\lambda>0$,
\[
\Prob\Bigl(\max_{1\le k\le n}\abs{S_k}\ge\lambda\Bigr)\le\frac{\Var(S_n)}{\lambda^2}.
\]
\end{theorem}
\begin{proof}
Let $A_k=\{\abs{S_k}\ge\lambda,\ \abs{S_j}<\lambda\ \text{for }j<k\}$, the event that
the first exceedance is at time $k$. The $A_k$ are disjoint and their union is the
event in question. Then
\[
\Exp[S_n^2]\ge\sum_{k=1}^n\Exp\bigl[S_n^2\indic_{A_k}\bigr]
\ge\sum_{k=1}^n\Exp\bigl[\bigl(S_k^2+2S_k(S_n-S_k)\bigr)\indic_{A_k}\bigr],
\]
using $S_n^2=S_k^2+2S_k(S_n-S_k)+(S_n-S_k)^2\ge S_k^2+2S_k(S_n-S_k)$. Now
$S_k\indic_{A_k}$ is a function of $X_1,\dots,X_k$ and $S_n-S_k$ of
$X_{k+1},\dots,X_n$, so they are independent and
$\Exp[S_k(S_n-S_k)\indic_{A_k}]=\Exp[S_k\indic_{A_k}]\Exp[S_n-S_k]=0$. Hence
$\Exp[S_n^2]\ge\sum_k\Exp[S_k^2\indic_{A_k}]\ge\lambda^2\sum_k\Prob(A_k)$, because
$S_k^2\ge\lambda^2$ on $A_k$.
\end{proof}

\begin{theorem}[Strong law of large numbers, finite-variance version]\label{thm:slln}
If $X_1,X_2,\dots$ are i.i.d.\ with mean $\mu$ and variance $\sigma^2<\infty$, then
$S_n/n\to\mu$ almost surely.
\end{theorem}
\begin{proof}
Replacing $X_i$ by $X_i-\mu$ we may assume $\mu=0$.

\emph{Step 1: along squares.} By Chebyshev, $\Prob(\abs{S_{k^2}}>\epsilon k^2)\le
k^2\sigma^2/(\epsilon^2k^4)=\sigma^2/(\epsilon^2k^2)$, which is summable in $k$. By
Theorem~\ref{thm:subseq}(i) (applied to the sequence $S_{k^2}/k^2$), $S_{k^2}/k^2\to0$
almost surely.

\emph{Step 2: between squares.} Let $D_k=\max_{k^2\le n<(k+1)^2}\abs{S_n-S_{k^2}}$.
The variables $S_n-S_{k^2}$ are partial sums of $X_{k^2+1},\dots$, a block of
$(k+1)^2-1-k^2=2k$ independent terms. By Kolmogorov's inequality,
\[
\Prob(D_k\ge\epsilon k^2)\le\frac{2k\sigma^2}{\epsilon^2k^4}=\frac{2\sigma^2}{\epsilon^2k^3},
\]
which is summable. By Theorem~\ref{thm:subseq}(i), $D_k/k^2\to0$ almost surely.

\emph{Step 3.} For $k^2\le n<(k+1)^2$ we have $n\ge k^2$ and
$\abs{S_n}\le\abs{S_{k^2}}+D_k$, so $\abs{S_n/n}\le\abs{S_{k^2}}/k^2+D_k/k^2\to0$
almost surely.
\end{proof}

\begin{nuance}{the strong law needs only a finite mean}
Kolmogorov's strong law: if $X_i$ are i.i.d.\ and $\Exp\abs{X_1}<\infty$, then
$S_n/n\to\Exp[X_1]$ almost surely; conversely if $\Exp\abs{X_1}=\infty$ then
$\limsup\abs{S_n}/n=\infty$ almost surely. The proof for a mere first moment uses
truncation and is due to Etemadi (see Durrett, \emph{Probability: Theory and
Examples}, Section~2.4); we use the result without proving it. The Cauchy
distribution (Example~\ref{ex:cauchycf}) shows that without a mean there is no law
of large numbers at all.
\end{nuance}

\begin{pitfall}{the law of large numbers does \emph{not} say that losses are ``due'' to be balanced}
The law says $S_n/n\to\mu$, not that $S_n-n\mu$ returns to zero. After a streak of
$10$ heads, the next toss is still $50\%$ heads. The excess heads are not cancelled
by later tails; they are diluted by the growing denominator. The ``gambler's
fallacy'' confuses these two statements.
\end{pitfall}

\begin{theorem}[Glivenko--Cantelli, continuous case]\label{thm:gc}
Let $X_1,X_2,\dots$ be i.i.d.\ with a continuous CDF $F$, and let
$F_n(x)=\frac1n\sum_{i=1}^n\indic_{\{X_i\le x\}}$ be the empirical CDF. Then
$\sup_x\abs{F_n(x)-F(x)}\to0$ almost surely.
\end{theorem}
\begin{proof}
Fix $m\ge1$ and pick points $-\infty=x_0<x_1<\dots<x_m=\infty$ with $F(x_j)=j/m$
(possible since $F$ is continuous). For each $j$, $F_n(x_j)\to F(x_j)$ almost surely by
the strong law applied to the indicators. So almost surely $F_n(x_j)\to F(x_j)$ for
all $j=0,\dots,m$ simultaneously (finite intersection). For $x\in[x_{j-1},x_j]$, by
monotonicity of $F_n$ and $F$,
\[
F_n(x)-F(x)\le F_n(x_j)-F(x_{j-1})=F_n(x_j)-F(x_j)+\tfrac1m,
\]
and similarly $F_n(x)-F(x)\ge F_n(x_{j-1})-F(x_{j-1})-\tfrac1m$. Hence
$\limsup_n\sup_x\abs{F_n(x)-F(x)}\le\tfrac1m$ almost surely. Let $m\to\infty$ along
the integers.
\end{proof}

\section{The central limit theorem}

The law of large numbers says $S_n/n-\mu\to0$. To see the fluctuations we rescale
by $\sqrt n$.

\begin{theorem}[Central limit theorem (Lindeberg--L\'evy)]\label{thm:clt}
Let $X_1,X_2,\dots$ be i.i.d.\ with mean $\mu$ and variance $0<\sigma^2<\infty$. Then
\[
\frac{S_n-n\mu}{\sigma\sqrt n}\xrightarrow{d}N(0,1).
\]
\end{theorem}

We need two lemmas.

\begin{lemma}[Taylor bound for $e^{ix}$]\label{lem:taylor}
For real $x$, define $R(x)=e^{ix}-1-ix+\tfrac12x^2$. Then
$\abs{R(x)}\le\min\bigl(\abs x^3/6,\ x^2\bigr)$.
\end{lemma}
\begin{proof}
Taylor's formula with integral remainder gives
$e^{ix}=1+ix+i^2\int_0^x(x-s)e^{is}ds$ and
$e^{ix}=1+ix-\tfrac12x^2+\frac{i^3}2\int_0^x(x-s)^2e^{is}ds$. Since $\abs{e^{is}}=1$,
the first gives $\abs{e^{ix}-1-ix}\le x^2/2$ and the second gives
$\abs{R(x)}\le\abs x^3/6$. Also
$\abs{R(x)}\le\abs{e^{ix}-1-ix}+x^2/2\le x^2$.
\end{proof}

\begin{lemma}[Products]\label{lem:products}
If complex numbers $a_1,\dots,a_n,b_1,\dots,b_n$ have modulus at most $1$, then
$\bigl\lvert\prod a_i-\prod b_i\bigr\rvert\le\sum\abs{a_i-b_i}$.
\end{lemma}
\begin{proof}
Telescope: $\prod_{i=1}^na_i-\prod_{i=1}^nb_i=\sum_{k=1}^n
\bigl(\prod_{i<k}b_i\bigr)(a_k-b_k)\bigl(\prod_{i>k}a_i\bigr)$, and each product of
unimodular-bounded factors has modulus at most $1$.
\end{proof}

\begin{proof}[Proof of Theorem~\ref{thm:clt}]
Replace $X_i$ by $(X_i-\mu)/\sigma$, so assume $\mu=0$, $\sigma=1$. Let $\varphi$ be the
characteristic function of $X_1$. Since the $X_i$ are independent,
\[
\varphi_{S_n/\sqrt n}(t)=\varphi\bigl(t/\sqrt n\bigr)^n .
\]
We must show this tends to $e^{-t^2/2}$ for each fixed $t$. Put $u=t/\sqrt n$. By
Lemma~\ref{lem:taylor} and $\Exp[X]=0$, $\Exp[X^2]=1$,
\[
\varphi(u)=\Exp[e^{iuX}]=1-\tfrac12u^2+\Exp[R(uX)]=1-\frac{t^2}{2n}+\Exp[R(uX)].
\]
Let $b=e^{-t^2/(2n)}$; then $b^n=e^{-t^2/2}$ exactly. Both $\abs{\varphi(u)}\le1$ and
$b\le1$, so by Lemma~\ref{lem:products}
\[
\bigl\lvert\varphi(u)^n-b^n\bigr\rvert\le n\abs{\varphi(u)-b}
\le n\Bigl\lvert1-\frac{t^2}{2n}-e^{-t^2/(2n)}\Bigr\rvert+n\,\Exp\abs{R(uX)}.
\]
For the first term use $\abs{1-z-e^{-z}}\le z^2/2$ for $z\ge0$ with $z=t^2/(2n)$: it is at
most $n\cdot t^4/(8n^2)\to0$. For the second,
\[
n\,\Exp\abs{R(uX)}\le\Exp\Bigl[\min\Bigl(\frac{\abs t^3\abs X^3}{6\sqrt n},\ t^2X^2\Bigr)\Bigr].
\]
For each outcome the quantity inside tends to $0$ as $n\to\infty$, and it is dominated
by $t^2X^2$, which is integrable. By dominated convergence the expectation tends to $0$.
Hence $\varphi_{S_n/\sqrt n}(t)\to e^{-t^2/2}$, the characteristic function of
$N(0,1)$ (Example~\ref{ex:cfnormal}). By L\'evy's continuity theorem
(Theorem~\ref{thm:cfprops}(vi), proved in Chapter~\ref{ch:measure}),
$S_n/\sqrt n\xrightarrow{d}N(0,1)$.
\end{proof}

\begin{corollary}[De Moivre--Laplace]
If $X_n\sim\mathrm{Bin}(n,p)$, then $(X_n-np)/\sqrt{npq}\xrightarrow{d}N(0,1)$.
\end{corollary}
\begin{proof}
$X_n$ is a sum of $n$ i.i.d.\ Bernoulli$(p)$ variables (mean $p$, variance $pq$).
\end{proof}

\begin{example}[Normal approximation with continuity correction]
For $X\sim\mathrm{Bin}(100,\tfrac12)$, $\mu=50$, $\sigma=5$. The event $\{X\ge60\}$
equals $\{X>59.5\}$ in the integer-valued sense, so
$\Prob(X\ge60)\approx1-\Phi\bigl(\tfrac{59.5-50}5\bigr)=1-\Phi(1.9)\approx0.0287$,
compared with the exact value $0.0284$. Without the $\tfrac12$ correction the answer
would be $1-\Phi(2)=0.0228$, noticeably worse.
\end{example}

\begin{nuance}{what the CLT does and does not say}
\begin{itemize}
\item The CLT is about the \emph{centre} of the distribution at scale $\sqrt n$. It
gives no information about tail probabilities at a fixed multiple of $n$ (use
Chernoff or Hoeffding), nor, for a fixed $n$, about the error. The
\emph{Berry--Esseen theorem} (cited, not proved) bounds the error:
$\sup_x\abs{F_n(x)-\Phi(x)}\le C\,\Exp\abs{X_1-\mu}^3/(\sigma^3\sqrt n)$ with
an absolute constant $C<0.5$. The rate $n^{-1/2}$ cannot be improved in general.
\item Finite variance is essential. For distributions with infinite variance the
normalised sums, if they converge at all, converge to \emph{stable laws} (the Cauchy
law is the stable law of index $1$), and the right scaling is not $\sqrt n$.
\item Convergence is in distribution: the sequence $(S_n-n\mu)/(\sigma\sqrt n)$ does
not converge almost surely, or in probability; it keeps fluctuating at order $1$
forever (this is made precise by the law of the iterated logarithm, cited).
\end{itemize}
\end{nuance}

The CLT extends to non-identically distributed summands under the \emph{Lindeberg
condition}, and to random vectors (the multivariate CLT, with limit
$N(0,\Sigma)$); we state these without proof.

\section{Slutsky's theorem and the delta method}

\begin{theorem}[Slutsky]\label{thm:slutsky}
Suppose $X_n\xrightarrow{d}X$ and $Y_n\xrightarrow{p}c$ for a constant $c$. Then
$X_n+Y_n\xrightarrow{d}X+c$ and $X_nY_n\xrightarrow{d}cX$.
\end{theorem}
\begin{proof}
\emph{Sum.} Let $x$ be a continuity point of $F_{X+c}$, i.e.\ $x-c$ is a continuity point of
$F_X$. For $\epsilon>0$ chosen so that $x-c\pm\epsilon$ are continuity points of $F_X$
(all but countably many $\epsilon$ qualify),
\[
\Prob(X_n\le x-c-\epsilon)-\Prob(\abs{Y_n-c}>\epsilon)\le\Prob(X_n+Y_n\le x)
\le\Prob(X_n\le x-c+\epsilon)+\Prob(\abs{Y_n-c}>\epsilon).
\]
Let $n\to\infty$: $F_X(x-c-\epsilon)\le\liminf\le\limsup\le F_X(x-c+\epsilon)$; now let
$\epsilon\downarrow0$ through admissible values, using continuity of $F_X$ at $x-c$.

\emph{Product.} Write $X_nY_n=cX_n+X_n(Y_n-c)$. First $cX_n\xrightarrow{d}cX$
(if $c\ne0$ the CDFs transform by rescaling the argument; if $c=0$ the limit is the point
mass at $0$ and this is trivial). Second, $X_n(Y_n-c)\xrightarrow{p}0$: given
$\delta,\eta>0$ choose $K$ so that $\Prob(\abs X\ge K)<\eta$ with $\pm K$ continuity
points of $F_X$; then
$\Prob(\abs{X_n(Y_n-c)}>\delta)\le\Prob(\abs{X_n}>K)+\Prob(\abs{Y_n-c}>\delta/K)$, whose
$\limsup$ is at most $\eta$. Conclude by the sum case applied to
$cX_n$ and $X_n(Y_n-c)$.
\end{proof}

\begin{theorem}[Delta method]\label{thm:delta}
Suppose $\sqrt n(T_n-\theta)\xrightarrow{d}N(0,\tau^2)$ and $g$ is differentiable at
$\theta$. Then $\sqrt n\bigl(g(T_n)-g(\theta)\bigr)\xrightarrow{d}N\bigl(0,\,g'(\theta)^2\tau^2\bigr)$.
\end{theorem}
\begin{proof}
First $T_n\xrightarrow{p}\theta$, since $T_n-\theta=\frac1{\sqrt n}\cdot\sqrt n(T_n-\theta)$
is a product of a sequence tending to $0$ and a convergent-in-distribution sequence
(Slutsky). By differentiability, $g(t)-g(\theta)=g'(\theta)(t-\theta)+r(t)(t-\theta)$
where $r(t)\to0$ as $t\to\theta$ ($r(\theta)=0$). Hence
\[
\sqrt n\bigl(g(T_n)-g(\theta)\bigr)=g'(\theta)\sqrt n(T_n-\theta)+r(T_n)\sqrt n(T_n-\theta).
\]
Since $T_n\xrightarrow{p}\theta$ and $r$ is continuous at $\theta$, $r(T_n)\xrightarrow{p}0$
(for any $\epsilon$ pick $\delta$ with $\abs{r(t)}<\epsilon$ for $\abs{t-\theta}<\delta$).
By Slutsky the last term tends to $0$ in probability and the first term tends in
distribution to $N(0,g'(\theta)^2\tau^2)$; apply Slutsky's sum rule.
\end{proof}

\begin{example}[Variance-stabilising transformation]
Let $\bar X_n$ be the average of $n$ i.i.d.\ Poisson$(\lambda)$ variables. By the CLT,
$\sqrt n(\bar X_n-\lambda)\xrightarrow{d}N(0,\lambda)$. With $g(x)=\sqrt x$,
$g'(\lambda)=1/(2\sqrt\lambda)$, so
$\sqrt n(\sqrt{\bar X_n}-\sqrt\lambda)\xrightarrow{d}N(0,\tfrac14)$. The limiting variance
does not depend on $\lambda$: the square root \emph{stabilises} the variance, which is
why $\sqrt{\text{count}}$ is used in practice.
\end{example}

\begin{interview}{which limit theorem?}
\emph{Averages converge}: LLN (weak if convergence in probability suffices, strong for
almost sure). \emph{Distribution of the average at scale $1/\sqrt n$}: CLT.
\emph{Probability of a large deviation at scale $n$}: Chernoff/Hoeffding/Cram\'er.
\emph{Smooth function of an average}: delta method. \emph{Infinite variance or
infinite mean}: the classical theorems do not apply; check for stable laws.
\end{interview}

\section*{Exercises}
\addcontentsline{toc}{section}{Exercises}

\begin{enumerate}
\item Prove that if $X_n\xrightarrow{p}X$ and $g$ is continuous then
$g(X_n)\xrightarrow{p}g(X)$ (first for uniformly continuous $g$, or for bounded
$X$).
\item Let $X_n\sim\mathrm{Exp}(n)$. Show $X_n\to0$ in probability, almost surely (use
Borel--Cantelli with independent copies), and in $L^1$.
\item Show that if $\sum_n\Prob(\abs{X_n}>\epsilon)<\infty$ for all $\epsilon$ then
$X_n\to0$ a.s.; give an example where $X_n\to0$ in probability but
$\sum\Prob(\abs{X_n}>\epsilon)=\infty$ for some $\epsilon$.
\item Let $X_1,X_2,\dots$ be i.i.d.\ $U(0,1)$ and $M_n=\max_{i\le n}X_i$. Show
$M_n\to1$ almost surely, and that $n(1-M_n)\xrightarrow{d}\mathrm{Exp}(1)$.
\item Use the CLT to approximate the probability that the sum of $100$ fair dice
exceeds $370$.
\item Show the CLT implies that $\Prob(S_n>n\mu)\to\tfrac12$. Does the
\emph{almost sure} version hold (does $S_n-n\mu$ change sign infinitely often)?
\item Prove that if $X_n\xrightarrow{d}X$ and $g$ is continuous then
$g(X_n)\xrightarrow{d}g(X)$ (continuous mapping theorem) for bounded continuous $g$,
using the characterisation $\Exp[g(X_n)]\to\Exp[g(X)]$ (cited).
\item Let $X_n$ be the number of fixed points of a random permutation of $n$ elements.
Show using inclusion--exclusion that $\Prob(X_n=k)\to e^{-1}/k!$ (Poisson$(1)$ limit).
\item (Challenge) Prove the weak law of large numbers for i.i.d.\ variables with
$\Exp\abs{X_1}<\infty$ only (truncate at level $\epsilon n$, use Chebyshev on the
truncated part and dominated convergence on the remainder).
\end{enumerate}
